%======================================================================
\section{Proof of the streaming theorem}\label{app:streaming}
%======================================================================

This appendix proves Theorem~\ref{thm:streaming}, the streaming service at the rank rate. Theorem~B uses it through
Corollary~\ref{cor:streaming-lamps}, which gives the services at the rank rate in Section~\ref{sec:main-channel-proof}.

The device keeps one rectangular dilation $W$ for all rounds and lets the part of the memory in use, the \emph{active space}, grow by the
factor $k'/k$ in every round; that growth is the whole purity bill. Three things are then controlled. The Kraus histories must stay nearly
orthogonal, and a Haar-random unitary of the active space before each round ensures this, through a concentration bound on the Gram matrix of
the histories. The concentration bound needs the state of everything outside the memory to stay spread out, and exporting about
$\frac12\log r$ qubits per round ensures that. And the error follows from the Gram bound alone, because the final vector of the service differs
from the ideal one by the polar part of a frame. The proof fixes the schedule, chooses the unitaries round by round, and then bounds the error
and the memory.

\begin{proof}[Proof of Theorem~\ref{thm:streaming}]
We use the frames of~\cite[Appendix~A]{DouglasMghirbiB}. A \emph{frame} is a linear map $F:\mathbb C^M\to\mathbb C^D\otimes\mathbb C^S$; its columns
are the vectors, in the memory $\mathbb C^D$ and an inaccessible exterior $\mathbb C^S$, attached to the Kraus histories. The exterior holds the
purifier of the initial memory, the exported qubits and the purifiers of the imported ones. Put $g(F)=\|F^*F\|$ and
$\alpha(F)=\|\mathrm{Tr}_{\mathbb C^D}FF^*\|$. Two lemmas from there carry the argument.
\begin{itemize}[leftmargin=1.2em,itemsep=1pt]
\item \emph{Concentration}~\cite[Lemma~A.1]{DouglasMghirbiB}. For Hermitian $X$ on $\mathbb C^D$ and Haar-random $V\in U(D)$, the matrix
$Z=F^*((V^*XV)\otimes I)F-\tau_D(X)F^*F$ satisfies $\Pr\{\|Z\|\ge s\}\le2M\exp[-Ds^2/(4\alpha(F)g(F)\|X\|^2)]$.
\item \emph{Export}~\cite[Lemma~A.4]{DouglasMghirbiB}. If $\mathbb C^D=\mathbb C^{D/e}\otimes\mathbb C^e$, there is a unitary of $\mathbb C^D$ after
which exporting the factor $\mathbb C^e$ and importing a maximally mixed $e$-dimensional register give a frame $F'$ with $F'^*F'=F^*F$ and, for every
$\eta>0$, $\alpha(F')\le(1+\eta)\alpha(F)/e^2+4e^2g(F)\ln(8Se^2)/(\eta D)$.
\end{itemize}

\emph{Fix the schedule and the dimensions.} Put $h=\log r$, $E_p=2^{\lceil h/2\rceil}$, $J_t=\lfloor th/2\rfloor$, $j_t=J_t-J_{t-1}\le\lceil h/2\rceil$ and
$e_t=2^{j_t}$, a divisor of $E_p$. Choose $D_0$, a multiple of $kE_p$, with $D_0\ge nkE_p$ and
\begin{equation}\label{eq:D0}
D_0^2\ \ge\ \frac{32\,A_*\,r^4n^2\ln(8r^{n+2})}{\epsilon^2},\qquad\text{where}\quad A_*=12\bigl(1+8E_p^2n^2L_*\bigr),\quad L_*=\ln(8D_0r^nE_p^2).
\end{equation}
Put $q_0=k'/k$. For $t\ge1$ let $\widetilde D_t=q_0D_{t-1}$, a multiple of $k'E_p$, and let $D_t$ be the least multiple of $kE_p$ with
$D_t\ge\widetilde D_t$. Then $D_t\le q_0D_{t-1}+kE_p$, so $D_n\le q_0^n(D_0+nkE_p)\le2q_0^nD_0$. The memory has $Q=\lceil\log D_n\rceil$ qubits, of
which the last $\lceil h/2\rceil$ form a \emph{port}. A space whose dimension $D'$ is a multiple of $E_p$ is realized as a subspace of the first
$Q-\lceil h/2\rceil$ qubits of dimension $D'/E_p$, tensored with the port. The active space of round $t$ is such a space of dimension $D_{t-1}$,
identified with $\mathbb C^k\otimes\mathbb C^{D_{t-1}/k}$.

\emph{The device never leaves the active spaces.} The memory starts maximally mixed on the active space of round $1$, of dimension $D_0$. In
round $t$ the device applies a unitary $V_t$ of the active space and then $W\otimes I_{D_{t-1}/k}$, which maps the input and the active space
isometrically onto the output and a space of dimension $\widetilde D_t$. It releases the output. It then applies a unitary $V'_t$ of that space;
since $V'_t$ acts on the memory alone, it could equally act before the release. If $t<n$, the device exports $j_t$ port qubits and imports
$j_t$ qubits in the maximally mixed state in their place. The resulting space of dimension $\widetilde D_t$ lies in the active space of round
$t+1$. These maps are completed arbitrarily to unitaries of the input and the whole memory; the completion never acts, since the memory never
leaves the active spaces. So $C=0$, $J=J_{n-1}$, and $P=Q-\log D_0\le n\log q_0+2$.

\emph{Frames track the histories.} Purify the initial memory: its frame $F_0$ is one unit vector, with $\alpha(F_0)=1/D_0$. Let $F$ be the frame
before round $t$, with one column $F_y$ for each history $y\in[r]^{t-1}$. After the call the frame has the columns $(B_i\otimes I)(V_t\otimes I)F_y$,
indexed by the histories $(y,i)$. Its Gram matrix has the blocks $F^*\bigl((V_t^*(B_i^*B_j\otimes I)V_t)\otimes I\bigr)F$, whose Haar means are
$\tau(B_i^*B_j\otimes I)F^*F=\delta_{ij}F^*F$ by the first identity in~\eqref{eq:rectangular}. By the second identity and cyclicity of the partial
trace over the memory, the exterior marginal after the call is exactly $r$ times the one before, whatever $V_t$ is; so the call multiplies
$\alpha$ by $r$.

\emph{Choose the unitaries round by round.} We choose $V_1,V'_1,V_2,\dots$ in turn so that the frame $F_t$ after round $t$ satisfies
\begin{equation}\label{eq:streaming-induction}
\|F_t^*F_t-I\|\le t\epsilon/n,\qquad \alpha(F_t)\le A_*/D_0 .
\end{equation}
This holds for $t=0$. Suppose it holds for $t-1$; then $g(F_{t-1})\le2$. Since $B_i^*B_i\le rI$, the Hermitian and anti-Hermitian parts of the $B_i^*B_j$
have norm at most $r$. Apply the concentration lemma to these $2r^2$ operators, tensored with $I_{D_{t-1}/k}$, with $D=D_{t-1}\ge D_0$, at most $r^n$
columns and level $s=\epsilon/(2rn)$: by~\eqref{eq:D0} the failure probabilities add up to at most
$4r^{n+2}\exp[-D_0^2\epsilon^2/(32A_*r^4n^2)]\le\frac12$. So some $V_t$ makes every block of the new Gram matrix deviate from its mean by at most
$\epsilon/(rn)$, hence the $r\times r$ block matrix by at most $\epsilon/n$, and the Gram matrix after the call is within $t\epsilon/n$ of $I$. Exports and
imports preserve it. Choose $V'_t$ by the export lemma with $e=e_t$, $\eta=1/n$ and $D=\widetilde D_t\ge D_0$. The exterior has dimension
$S\le D_0\,2^{2J_{t-1}}\le D_0r^n$, so $\ln(8Se_t^2)\le L_*$, and
\[
\alpha(F_t)\ \le\ \Bigl(1+\frac1n\Bigr)\frac{r\,\alpha(F_{t-1})}{e_t^2}+\frac{8E_p^2nL_*}{D_0}.
\]
Put $c_t=2^{2J_t}r^{-t}$, so that $c_0=1$, $\frac14<c_t\le1$ and $c_tr/e_t^2=c_{t-1}$. Multiplying by $c_t$ and unrolling gives
$c_t\alpha(F_t)\le(1+\frac1n)^t(1+8E_p^2n^2L_*)/D_0$, hence $\alpha(F_t)\le4e(1+8E_p^2n^2L_*)/D_0\le A_*/D_0$. In round $n$ there is no export, and only the
Gram bound is needed.

\emph{The error is at most $\epsilon$.} Purify the observer. In the ideal experiment, with the Stinespring isometry
$\xi\mapsto\sum_iA_i\xi\otimes|i\rangle$ in each round, its final vector is $\sum_y\psi_y\otimes|y\rangle$. In the service it is
$\sum_y\psi_y\otimes F_n|y\rangle$ with the same vectors $\psi_y$, because the device acts on the input and the memory through
$W=\sum_iA_i\otimes B_i$ and unitaries of the memory, and the observer never touches the memory or the exterior. The polar part
$T=F_n(F_n^*F_n)^{-1/2}$ is an isometry with $\|F_n-T\|=\|(F_n^*F_n)^{1/2}-I\|\le\epsilon$. So the actual final vector is within $\epsilon$ of the
ideal one mapped by $I\otimes T$, and after tracing out the memory and the exterior the observer's final states are within $\epsilon$ in trace
distance.

\emph{The memory is as claimed.} We have $Q\le\log D_n+1\le\log D_0+n\log q_0+2$. The conditions on $D_0$ hold for
$D_0=kE_p\lceil K_r(n/\epsilon)^4(1+\ln k)\rceil$ with $K_r$ large enough in terms of $r$, and then
$\log D_0=\log k+O_r(\log(n/\epsilon)+\log(2+\log k))$.
\end{proof}
